% !TEX program = xelatex
\documentclass[11pt,openany,fontset=fandol]{ctexbook}
\usepackage[paperwidth=176mm,paperheight=250mm,inner=21mm,outer=18mm,top=20mm,bottom=20mm,headheight=14pt,headsep=6mm,footskip=10mm]{geometry}
\usepackage{amsmath,amssymb,amsthm,bm,mathtools,booktabs,array,enumitem,float}
\usepackage{tikz,fancyhdr,hyperref,graphicx}
\usetikzlibrary{arrows.meta,calc,patterns,angles,quotes,decorations.markings,decorations.pathmorphing}
\theoremstyle{definition}\newtheorem{example}{例}[chapter]\newtheorem{drill}{独立练习}[chapter]
\newenvironment{solution}{\par\smallskip\noindent\textbf{解}\quad\ignorespaces}{\par\medskip}
\newenvironment{drillsolution}[1]{\subsubsection*{独立练习\ref{#1}详解}\phantomsection\label{sol:#1}\ignorespaces}{\par\medskip}
\newenvironment{question}{\par\medskip\noindent\begin{minipage}{\linewidth}}{\end{minipage}\par\medskip}
\newcommand{\dd}{\mathop{}\!\mathrm d}
\newcommand{\e}{\mathrm e}
\newsavebox{\AstroFigureBox}\newcommand{\fig}[3]{\begin{figure}[H]\centering\sbox{\AstroFigureBox}{#3}\ifdim\wd\AstroFigureBox>\linewidth\resizebox{\linewidth}{!}{\usebox{\AstroFigureBox}}\else\usebox{\AstroFigureBox}\fi\caption{#2}\label{#1}\end{figure}}
\newcommand{\hintlink}[1]{\par\smallskip\noindent\hyperref[hint:#1]{提示}\quad\hyperref[sol:#1]{详解}}
\tikzset{axis/.style={->,>=Stealth,line width=.45pt},world/.style={line width=.8pt},flow/.style={->,>=Stealth,line width=.7pt},ray/.style={->,>=Stealth,line width=.7pt},guide/.style={dashed,line width=.35pt},aux/.style={dashed,line width=.35pt},dim/.style={<->,>=Stealth,line width=.35pt},every node/.style={font=\small},force/.style={->,>=Stealth,very thick},gas/.style={fill=blue!7}}
% Common visual style for all first-round physics lecture projects.
% Each project includes an identical local copy of this file.
\geometry{paperwidth=176mm,paperheight=250mm,inner=21mm,outer=18mm,top=20mm,bottom=20mm,headheight=14pt,headsep=6mm,footskip=10mm}
\hypersetup{hidelinks}
\setlength{\parindent}{2em}\setlength{\parskip}{3pt}\setlength{\emergencystretch}{2em}
\linespread{1.18}\raggedbottom\setcounter{tocdepth}{1}
\widowpenalty=10000\clubpenalty=10000\displaywidowpenalty=10000
\setlength{\intextsep}{8pt plus1pt minus1pt}
\makeatletter
\g@addto@macro\normalsize{\setlength{\abovedisplayskip}{8pt plus2pt minus2pt}\setlength{\belowdisplayskip}{8pt plus2pt minus2pt}\setlength{\abovedisplayshortskip}{0pt plus2pt}\setlength{\belowdisplayshortskip}{6pt plus2pt minus2pt}}
\makeatother
\providecommand{\Needspace}[1]{}
\newdimen\LectureAvailableSpace
\renewcommand{\Needspace}[1]{\par\penalty-100\begingroup\LectureAvailableSpace=\pagegoal\advance\LectureAvailableSpace by-\pagetotal\ifdim\LectureAvailableSpace<#1\relax\ifdim\pagetotal>0pt\newpage\fi\fi\endgroup}
\renewenvironment{solution}{\par\smallskip\Needspace{3\baselineskip}\noindent\textbf{解}\quad\ignorespaces}{\par\medskip}
\ctexset{chapter={name={第,讲},number=\chinese{chapter},format=\LARGE\bfseries\centering,beforeskip=16pt,afterskip=22pt},section={format=\Large\bfseries},subsection={format=\large\bfseries},subsubsection={format=\normalsize\bfseries,beforeskip=10pt,afterskip=4pt}}
\setlist[enumerate]{leftmargin=2em,itemsep=2pt,topsep=3pt}
\pagestyle{fancy}\fancyhf{}\fancyhead[LE,RO]{\small\thepage}\fancyhead[LO]{\small 物理竞赛一轮讲义}
\renewcommand{\headrulewidth}{0pt}
\fancypagestyle{plain}{\fancyhf{}\fancyfoot[C]{\small\thepage}\renewcommand{\headrulewidth}{0pt}}
\renewenvironment{drillsolution}[1]{\par\Needspace{4\baselineskip}\subsubsection*{独立练习\ref{#1}详解}\phantomsection\label{sol:#1}\nobreak\noindent\textbf{解}\quad\ignorespaces}{\par\medskip}
\ifdefined\exercisesolution
\renewenvironment{exercisesolution}[1]{\par\Needspace{4\baselineskip}\subsubsection*{习题\ref{#1}详解}\phantomsection\label{sol:#1}\nobreak\noindent\textbf{解}\quad\ignorespaces}{\par\medskip}
\fi
\newcommand{\LectureAnswerSection}{\par\section*{练习提示与解答}\addcontentsline{toc}{section}{练习提示与解答}}
\newcommand{\LectureHint}[3]{\par\Needspace{4\baselineskip}\paragraph{练习\ref{#1}}\noindent\textbf{提示一：}#2\par\noindent\textbf{提示二：}#3\par}
\newcommand{\LectureHeader}[2]{\fancyhead[RE]{\small #1\quad #2}}
\newcommand{\LectureCover}[3]{\begin{titlepage}\centering\vspace*{32mm}{\Huge\bfseries #1\par}\vspace{12mm}{\Large 物理竞赛一轮讲义\par}\vspace{18mm}{\LARGE #2\\[5mm]#3\par}\vfill{\large 2026\par}\end{titlepage}}

\LectureHeader{第四讲}{逃逸、碰撞与引力助推}
\begin{document}\frontmatter\LectureCover{天体运动}{第四讲}{逃逸、碰撞与引力助推}\tableofcontents\mainmatter
\setcounter{chapter}{3}\chapter{逃逸、碰撞与引力助推}
\section{逃逸是相对运动的边界条件}
\subsection{无穷远剩余速度与表面速度}

在固定球形天体外，单位质量能量 $\varepsilon=v^2/2-\mu/r$。若轨道可向外延伸到无穷远，则 $\varepsilon=v_\infty^2/2$，所以局部速率满足 $v^2=v_\infty^2+2\mu/r$。$v_\infty$ 称双曲余速；逃逸的临界值是 $v_\infty=0$，对应当地逃逸速度 $v_{\rm esc}=\sqrt{2\mu/r}$。

能量非负只说明轨道存在延伸到无穷远的分支。若当前速度径向向内且角动量为零，物体会先撞上中心，不能把“存在数学上的出射分支”当作实际逃逸。对有尺寸的天体，还要检查近地点是否在表面之外；对能反弹、能穿透或会被吸收的物体，后续边界也不同。

两体都能运动时，在质心系使用相对速度 $\bm u=\bm v_1-\bm v_2$，相对引力参数为 $G(m_1+m_2)$。两体分离到无穷远的能量阈值由相对运动决定，整体质心在实验室中平动得再快，也不会单独帮助两体摆脱相互束缚。

\subsection{相落时间的正、零、负能量分支}

纯径向向内的初速为 $u_0$、初间距 $l$，能量式给 $\dot r=-\sqrt{u_0^2+2\mu(1/r-1/l)}$。相遇时间应从此负根换成正的积分
\[t=\int_0^l\frac{\dd r}{\sqrt{u_0^2+2\mu(1/r-1/l)}}.\]
这个表达式对三种能量都成立。若令 $k=2\mu/l-u_0^2$，$k>0$ 用三角换元，$k=0$ 直接幂积分，$k<0$ 用双曲函数换元。只写负能量的反三角函数而不检验初速范围，会在高速入射时得到非实数时间。


\section{双曲线、散射与碰撞截面}
\subsection{几何参数的正负约定}

标准双曲线 $x^2/a_h^2-y^2/b_h^2=1$ 中 $a_h,b_h>0$，$c_h^2=a_h^2+b_h^2$、$e=c_h/a_h>1$。渐近线倾角满足 $\tan\theta=b_h/a_h$，焦点到渐近线的垂距恰为 $b_h$。这个垂距是无穷远入射的冲量参数，并非近地点距离。

引力双曲线能量为正：$\varepsilon=\mu/(2a_h)$，而角动量满足 $h^2=\mu a_h(e^2-1)=\mu b_h^2/a_h$。若沿用活力公式的有符号半长轴，则应取 $a=-a_h<0$，不能把正的几何 $a_h$ 直接塞进椭圆的负能量式。

\typeout{ASTROFIG:fig:at4:hyperbola}\fig{fig:at4:hyperbola}{入射渐近线与最近点属于不同几何量}{\begin{tikzpicture}[scale=.85]
\draw[guide](-4.6,-1.118)--(1,-1.118);\draw[world,domain=-2:2,samples=110]plot({(1.5-cosh(\x))/1.5+1.11803399*1.11803399*sinh(\x)/1.5},{-1.11803399*(1.5-cosh(\x))/1.5+1.11803399*sinh(\x)/1.5});\fill(0,0)circle(2pt)node[above left]{$F$};\draw[flow](-4.2,-1.118)--(-3,-1.118)node[midway,below]{$v_\infty$};\draw[dim](1,-1.118)--(1,0)node[midway,right]{$b$};\draw[guide](0,0)--(.333,-.373);\node[right]at(.5,-.25){$r_p$};
\end{tikzpicture}}
图\ref{fig:at4:hyperbola}中的 $b$ 垂直于入射渐近线，而 $r_p$ 从焦点指向最近点。两者之差反映了引力聚焦，图上不可把这两段长度混用。\subsection{散射角与最近距离}

无穷远入射速度 $v_\infty$、冲量参数 $b$，给 $h=bv_\infty$。由能量与离心率式，
\[e=\sqrt{1+\frac{b^2v_\infty^4}{\mu^2}},\quad
r_p=\frac{h^2/\mu}{1+e}=\frac\mu{v_\infty^2}(e-1).\]
吸引散射的速度转向角满足 $\sin(\chi/2)=1/e$，等价于 $\tan(\chi/2)=\mu/(bv_\infty^2)$。这里比较的是入射与出射速度，而不是两条朝外径矢之间的张角。远处速率不变，方向改变；若两端还保留同一中心势，能量守恒不等于速度矢量守恒。

对于半径 $R$ 的吸收体，刚好擦到表面的轨道有 $r_p=R$。消去近点速度或离心率，得到最大冲量参数 $b_{\max}^2=R^2(1+2\mu/(Rv_\infty^2))$。入射截面是 $\pi b_{\max}^2$，通常大于几何截面 $\pi R^2$。该公式假定稀薄粒子、无相互遮挡、无阻力以及命中后确实被吸收。


\section{碰撞改变轨道时分两段计算}
\subsection{短时动量守恒与随后机械能守恒}

陨石与卫星粘在一起，碰撞时间很短，可忽略这段时间中心引力的冲量，因此两物体的总动量近似守恒。非弹性碰撞一般损失机械能，不能跨碰撞全过程写卫星机械能守恒。求出碰后速度矢量后，再把合体看作新的试验物体，计算后续的 $\varepsilon,h,r_p$。

即使陨石沿径向撞入，卫星原有的切向动量仍在。碰后切向速度减小，但一般不为零，因此不能只比较碰后总能量与表面势能；必须给表面剩余切向动能留出能量。以到达表面时径向速度平方非负为条件，常比先求完整椭圆参数更直接。

\subsection{接触等号和时间方向}

若题目允许擦碰，边界是表面处径向速度为零；要求进入或撞击具有有限法向速度，则需要严格不等式。得到一个阈值表达式以后，还要检查它是否为正：如果无径向陨石速度时，单是切向减速就已使轨道进入表面，所求最低径向速度就是零，而不是对一个负数强行开根。


\section{移动行星的引力助推}
\subsection{先在行星系散射，再回到恒星系}

在局部掠近期间，近似行星的恒星系速度 $\bm V$ 不变，忽略恒星在行星附近的差分场。入、出行星引力影响区时，先作速度变换 $\bm u_{\rm in}=\bm v_{\rm in}-\bm V$。行星系双曲散射保持余速大小 $u$，只把方向旋转一个允许角 $\delta$，然后变回 $\bm v_{\rm out}=\bm V+\bm u_{\rm out}$。

在这个近似中，恒星系单位质量能量变化为 $\bm V\cdot(\bm u_{\rm out}-\bm u_{\rm in})$。它可以正也可以负，取决于从行星的哪一侧绕过。助推是与移动天体交换动量和能量，不是静止引力场无代价增加同一惯性系中的机械能。

\typeout{ASTROFIG:fig:at4:frames}\fig{fig:at4:frames}{恒星系速度是行星平动与相对速度的矢量和}{\begin{tikzpicture}[scale=.85]
\draw[flow](0,0)--(2.5,0)node[midway,below]{$\bm V$};\draw[flow](2.5,0)--(3.5,1.8)node[midway,right]{$\bm u$};\draw[force](0,0)--(3.5,1.8)node[midway,above left]{$\bm v$};
\end{tikzpicture}}
图\ref{fig:at4:frames}的速度三角形适用于入射和出射两端。只给速度大小而不给绕行一侧和冲量参数，通常不能得到唯一的出射矢量。\subsection{最大转角不能任意取反向}

若行星半径、大气或安全条件规定近行星点 $r_p\ge r_{\min}$，对给定余速 $u$，允许转角不超过
\[\delta_{\max}=2\arcsin\frac1{1+r_{\min}u^2/\mu_p}.\]
有限正的 $r_{\min},u$ 给 $\delta_{\max}<\pi$。把相对速度无条件反转180度，会夸大助推能力，有时对应穿入行星或极限冲量参数为零的不可行路径。

从地表逃离行星、沿恒星轨道转移、再掠近另一行星，是几个不同的保守飞行段。先用地球系求脱离余速，再变到日心系；到火星附近再转到火星系；离开火星后才用日心能量判断逃逸太阳。最小发射速度还依赖发射方向、会合相位和最小掠近半径，不能只把两个逃逸速度标量相加。


\section{散射中的时间、捕获与动量交换}
\subsection{双曲近点前后的时间}

  正能量轨道可用双曲偏近点参数 $H$ 表示，取近点为 $H=0$，$r=a_h(e\cosh H-1)$。面积速度不变给 $t-t_p=\sqrt{a_h^3/\mu}(e\sinh H-H)$。其导数因子 $e\cosh H-1>0$，所以同一有向分支上的时间随 $H$ 单调增加。

  初位置在入射支，参数应取负值；到近点的剩余时间取正。把余弦双曲函数反函数得到的正值直接当成初始 $H$，会把入射位置误认成出射位置。与椭圆不同，双曲线没有周期，也不能用某个“完整双曲线面积”除以面积速度求飞行时间。
  
\subsection{静止天体的保守引力不能自行捕获}

  无穷远以非零余速入射的物体，相对能量为正。若只有固定中心的保守引力，它在远处出射时仍有同样的正能量；单纯靠轨道弯曲不能把双曲线变成椭圆。捕获需要阻力、碰撞、点火或第三个天体交换能量，或者改变了所选系统边界。

  若在近点逆行点火，局部速率为 $\sqrt{v_\infty^2+2\mu/r_p}$。减到当地抛物线速率 $\sqrt{2\mu/r_p}$ 只是零能量边界，要形成真正椭圆还需稍多减速。取最小近点能减少所需速度增量，但同时受到实体半径、大气、潮汐和设备约束。
  
\subsection{吸积物质带来的动量变化}

  天体吸收静止尘埃时，合并前后动量关系为 $M\bm V=(M+\dd M)(\bm V+\dd\bm V)$，保留一阶得 $M\dd\bm V=-\bm V\dd M$。因此仅就吸收静止物质而言，天体会减速；其动能减少量并不全由平动速度变化项表达，还应包括质量增加项。

  引力聚焦只决定哪些远处粒子能命中，并不取消它们原来的动量条件。要计算长期轨道演化，必须把粒子原来的日心速度、吸收过程与天体受外力一并纳入，不能把“吸积率为正”直接等同于“天体获得推力加速”。
  

\section{分段二体近似与完整可行性}
\subsection{同一个无穷远在不同近似中的含义}

  掠近行星时常把行星附近看成二体系统，把足够远处的相对运动看成直线渐近线。这个“行星系无穷远”是局部二体轨道外推得到的余速，不表示在真实太阳系中可以一路忽略太阳，直到空间位置真的趋于无穷。离开局部区域后，应把渐近余速与行星公转速度相加，接回日心运动。

  检验局部行星主导，应比较太阳对探测器和行星的差分加速度，而非太阳对二者共同的加速度。若局部尺度为 $\ell$、行星日心距离为 $r_p$，太阳潮汐量级为 $\mu_s\ell/r_p^3$，行星引力为 $\mu_p/\ell^2$。前者远小于后者要求
  \[\ell\ll r_p(\mu_p/\mu_s)^{1/3}.\]
  这是量级条件，不在此承诺精确的影响区边界。遇到非常慢、非常远或长期绕行的掠近，太阳扰动和行星曲线运动都可能积累到不可忽略，局部固定速度散射模型应重新检查。
  
\typeout{ASTROFIG:fig:at4:patch}\fig{fig:at4:patch}{局部行星散射与日心飞行在共同边界上接续}{\begin{tikzpicture}[scale=.85]
\draw[guide](0,0)circle(1.6);\fill(0,0)circle(4pt)node[below]{行星};\draw[flow](-3,-1)--(-1.4,-.8)node[midway,above]{入界};\draw[world](-1.4,-.8)..controls(0,-1)and(1.2,-.5)..(1.4,.8);\draw[flow](1.4,.8)--(2.8,1.9)node[midway,left]{出界};\node at(0,-2.2){边界上转换速度参考系};
\end{tikzpicture}}

  图\ref{fig:at4:patch}的虚线圈是近似区域示意，并非一个有力的实体边界。边界两侧的位置和惯性系速度必须连续；变化的是计算中保留哪一个主导天体。不能在入界时把相对速度直接替换成日心速度，或在出界时把行星速度忘记加回。

  条件不足的助推题可以先列出可确定量和剩余自由度。能量式确定余速大小，角动量或冲量参数确定转角，绕行一侧确定转角符号，行星公转速度则负责把结果变回日心系。缺少其中任一项，可能只能给速度范围或一个参数族。保留参数不是省略计算，而是防止凭空增加题设没有给出的条件。

  有限安全半径还会把可行速度范围缩小。对同一个余速，转角有上限；对同一个目标出射方向，还要问这个方向是否在允许转向扇区内。先在速度图中画出可转动的扇区，再最大化日心速度投影，比直接设定“末速度必与行星速度同向”更稳妥。若那个方向不在扇区内，最优值在扇区边界而不是理想平行方向。

  实际最值的结论应区分已取得的最小值、只可逼近的下确界，以及由必要条件得到的下界。抛物线捕获边界、无穷中间半径和极限180度反转分别说明了三种容易混淆的情形。最后回代检验表面、方向、相位和时间，是让代数结果成为物理解答的必要一步。
  

\section{例题与分析}

\begin{question}\begin{example}[双曲线几何与引力守恒量]\label{at4:example:1}
双曲线 $x^2/a_h^2-y^2/b_h^2=1$，$a_h,b_h>0$，焦距 $c_h=\sqrt{a_h^2+b_h^2}$。求渐近线倾角、焦点到渐近线距离；若它是引力参数 $\mu$ 的轨道，求单位质量能量、角动量和近点距离。
\end{example}
\typeout{ASTROFIG:fig:at4:example:1:s}\fig{fig:at4:example:1:s}{题设：双曲线几何与引力守恒量}{\begin{tikzpicture}[scale=.85]
\draw[guide](-3,-2)--(3,2)(-3,2)--(3,-2);\draw[world,domain=-1.2:1.2,samples=80]plot({cosh(\x)},{.666667*sinh(\x)});\fill(1.20185,0)circle(2pt)node[below]{$F$};\draw[axis](-2.5,0)--(3.5,0);
\end{tikzpicture}}
\end{question}
\begin{solution}
图\ref{fig:at4:example:1:s} 的渐近线为 $y=\pm(b_h/a_h)x$，故 $\theta=\arctan(b_h/a_h)$。焦点 $(c_h,0)$ 到其中一条直线的距离
\[d=\frac{(b_h/a_h)c_h}{\sqrt{1+b_h^2/a_h^2}}=b_h.\]
\typeout{ASTROFIG:fig:at4:example:1:p}\fig{fig:at4:example:1:p}{解答过程：双曲线几何与引力守恒量}{\begin{tikzpicture}[scale=.85]
\draw[guide](0,0)--(3,2);\fill(2,0)circle(2pt)node[below]{$F$};\draw[dim](2,0)--(1.3846,.9231)node[midway,right]{$d=b_h$};\draw[axis](-.5,0)--(3.5,0);\node at(1,1.7){$\tan\theta=b_h/a_h$};
\end{tikzpicture}}

图\ref{fig:at4:example:1:p} 取 $e=c_h/a_h$，近点距离 $r_p=c_h-a_h=a_h(e-1)$。轨道半通径为 $p=a_h(e^2-1)=b_h^2/a_h$，所以
\[h=\sqrt{\frac{\mu b_h^2}{a_h}},\qquad \varepsilon=\frac\mu{2a_h}>0.\]
若试验物体质量为 $m$，总能量和角动量分别再乘 $m$。由 $v_\infty^2=2\varepsilon$ 得 $h=b_hv_\infty$，也与 $d=b_h$ 的几何意义一致。几何半长轴为正，不能在此写成椭圆的负能量。\end{solution}
\begin{question}\begin{example}[给定近地点的引力散射角]\label{at4:example:2}
小天体以无穷远速度 $v_\infty$ 掠过质量 $M$ 的固定天体，最近距离为 $r_p$。求偏转角和冲量参数，忽略中心天体尺寸。
\end{example}
\typeout{ASTROFIG:fig:at4:example:2:s}\fig{fig:at4:example:2:s}{题设：给定近地点的引力散射角}{\begin{tikzpicture}[scale=.85]
\draw[guide](-4.6,-1.118)--(1,-1.118);\draw[world,domain=-2:2,samples=110]plot({(1.5-cosh(\x))/1.5+1.11803399*1.11803399*sinh(\x)/1.5},{-1.11803399*(1.5-cosh(\x))/1.5+1.11803399*sinh(\x)/1.5});\fill(0,0)circle(2pt)node[above left]{$F$};\draw[flow](-4.2,-1.118)--(-3,-1.118)node[midway,below]{$v_\infty$};
\end{tikzpicture}}
\end{question}
\begin{solution}
图\ref{fig:at4:example:2:s} 无穷远能量为 $v_\infty^2/2$，因此几何半长轴 $a_h=\mu/v_\infty^2$。由 $r_p=a_h(e-1)$，
\[e=1+\frac{r_pv_\infty^2}{\mu}.\]
\typeout{ASTROFIG:fig:at4:example:2:p}\fig{fig:at4:example:2:p}{解答过程：给定近地点的引力散射角}{\begin{tikzpicture}[scale=.85]
\draw[guide](-4.6,-1.118)--(1,-1.118);\draw[world,domain=-2:2,samples=110]plot({(1.5-cosh(\x))/1.5+1.11803399*1.11803399*sinh(\x)/1.5},{-1.11803399*(1.5-cosh(\x))/1.5+1.11803399*sinh(\x)/1.5});\fill(0,0)circle(2pt)node[above left]{$F$};\draw[flow](-4.2,-1.118)--(-3,-1.118)node[midway,below]{$v_\infty$};\draw[guide](0,0)--(.333,-.373)node[midway,above right]{$r_p$};\node at(-2,2){$e=1+r_pv_\infty^2/\mu$};
\end{tikzpicture}}

图\ref{fig:at4:example:2:p} 两渐近线给速度转角 $\chi=2\arcsin(1/e)$，所以
\[\chi=2\arcsin\frac1{1+r_pv_\infty^2/\mu},\qquad
b^2=r_p^2+\frac{2\mu r_p}{v_\infty^2}.
\]
第二式也可由 $h^2=r_p^2(v_\infty^2+2\mu/r_p)=b^2v_\infty^2$ 检查。弱引力极限 $\mu\to0$ 时 $b\to r_p,\chi\to0$；固定 $r_p$ 而减小入射速度时转角增大。这里讨论无碰撞质点轨道，不能在有半径的天体中无限减小 $r_p$。\end{solution}
\begin{question}\begin{example}[有冲量参数的碰撞条件]\label{at4:example:3}
半径 $R$、质量 $M$ 的球形吸收体静止。小粒子在无穷远速度为 $v_\infty>0$，入射直线到中心距离为 $b$。求能够碰到表面的条件，允许相切。
\end{example}
\typeout{ASTROFIG:fig:at4:example:3:s}\fig{fig:at4:example:3:s}{题设：有冲量参数的碰撞条件}{\begin{tikzpicture}[scale=.85]
\draw[guide](-4.6,-1.118)--(1,-1.118);\draw[world,domain=-2:2,samples=110]plot({(1.5-cosh(\x))/1.5+1.11803399*1.11803399*sinh(\x)/1.5},{-1.11803399*(1.5-cosh(\x))/1.5+1.11803399*sinh(\x)/1.5});\fill(0,0)circle(2pt)node[above left]{$F$};\draw[flow](-4.2,-1.118)--(-3,-1.118)node[midway,below]{$v_\infty$};\draw[world](0,0)circle(.5);\node at(1.8,.2){表面 $R$};
\end{tikzpicture}}
\end{question}
\begin{solution}
图\ref{fig:at4:example:3:s} 用无穷远初态给 $h=bv_\infty$、$\varepsilon=v_\infty^2/2$。在球面若可达，径向速度平方须非负：
\[v_r(R)^2=v_\infty^2+\frac{2\mu}{R}-\frac{b^2v_\infty^2}{R^2}\ge0.
\]
\typeout{ASTROFIG:fig:at4:example:3:p}\fig{fig:at4:example:3:p}{解答过程：有冲量参数的碰撞条件}{\begin{tikzpicture}[scale=.85]
\draw[guide](-4.6,-1.118)--(1,-1.118);\draw[world,domain=-2:2,samples=110]plot({(1.5-cosh(\x))/1.5+1.11803399*1.11803399*sinh(\x)/1.5},{-1.11803399*(1.5-cosh(\x))/1.5+1.11803399*sinh(\x)/1.5});\fill(0,0)circle(2pt)node[above left]{$F$};\draw[flow](-4.2,-1.118)--(-3,-1.118)node[midway,below]{$v_\infty$};\draw[world](0,0)circle(.5);\draw[dim](1,-1.118)--(1,0)node[midway,right]{$b$};\node at(-2,2){$r_p=R$ 为临界};
\end{tikzpicture}}

图\ref{fig:at4:example:3:p} 的相切边界给 $b_{\max}^2=R^2+2\mu R/v_\infty^2$。若 $b\le R$，所有正的 $v_\infty$ 都可命中；若 $b>R$，则
\[v_\infty^2\le\frac{2\mu R}{b^2-R^2}.\]
方向是不等号向下：速度越快，引力偏转越小，过大的冲量参数来不及被聚焦到表面。等号是擦碰；若必须穿入表面，取严格小于。这里的判断还利用了轨道从无穷远沿入射支到近点，不能拿它直接处理初始已经向外的轨道。\end{solution}
\begin{question}\begin{example}[行星扫过尘埃的吸积率]\label{at4:example:4}
半径 $R$、质量 $M$ 的行星以速度 $V$ 穿过静止、均匀、密度为 $\rho$ 的稀薄尘埃。尘埃碰到表面即被吸收，忽略粒子间作用和行星参数的瞬时变化，求吸积质量率。
\end{example}
\typeout{ASTROFIG:fig:at4:example:4:s}\fig{fig:at4:example:4:s}{题设：行星扫过尘埃的吸积率}{\begin{tikzpicture}[scale=.85]
\draw[world](0,0)circle(1);\draw[flow](0,0)--(2,0)node[above]{$V$};\foreach \x in{-3,-2,3,4}{\foreach \y in{-1.5,-.5,.5,1.5}{\fill(\x,\y)circle(.8pt);}}
\end{tikzpicture}}
\end{question}
\begin{solution}
图\ref{fig:at4:example:4:s} 换到行星系，远处尘埃以 $V$ 入射。上一碰撞条件给 $b_{\max}^2=R^2(1+2GM/(RV^2))$。
\typeout{ASTROFIG:fig:at4:example:4:p}\fig{fig:at4:example:4:p}{解答过程：行星扫过尘埃的吸积率}{\begin{tikzpicture}[scale=.85]
\draw[world](0,-1.5)rectangle(4,1.5);\draw[dim](-.5,-1.5)--(-.5,1.5)node[midway,left]{$2b_{\max}$};\draw[dim](0,-2)--(4,-2)node[midway,below]{$V\dd t$};\node at(2,0){入射质量 $\rho\pi b_{\max}^2V\dd t$};
\end{tikzpicture}}

图\ref{fig:at4:example:4:p} 在时间 $\dd t$ 内被捕获的是远处半径 $b_{\max}$、长度 $V\dd t$ 的柱体质量，故
\[\frac{\dd M}{\dd t}=\rho\pi R^2V\left(1+\frac{2GM}{RV^2}\right).
\]
括号内第一项是几何扫掠，第二项是引力聚焦。单位为质量每时间，符合质量流量。形式上 $V\to0$ 时聚焦项发散，提示匀速、无限均匀稀薄流和瞬时稳态近似不能任意延伸到零速；真实环境的有限尺度和其他相互作用会限制它。\end{solution}
\begin{question}\begin{example}[陨石粘合后卫星是否撞地]\label{at4:example:5}
质量 $m$ 的卫星绕质量 $M$、半径 $R_0$ 的行星作半径 $R>R_0$ 的圆运动。质量 $m/10$ 的陨石沿径向向内、以速度 $u$ 撞入并粘合。求能使合体碰到表面的最低 $u$，允许擦碰。
\end{example}
\typeout{ASTROFIG:fig:at4:example:5:s}\fig{fig:at4:example:5:s}{题设：陨石粘合后卫星是否撞地}{\begin{tikzpicture}[scale=.85]
\draw[world](0,0)circle(1);\draw[guide](0,0)circle(2);\fill(2,0)circle(2pt)node[below]{$m$};\draw[flow](2,0)--(2,1)node[right]{$v_c$};\draw[force](3.5,0)--(2.6,0)node[above]{$u,\ m/10$};\draw[dim](0,-.5)--(2,-.5)node[midway,below]{$R$};\draw[dim](0,0)--(-.707,.707)node[midway,left]{$R_0$};
\end{tikzpicture}}
\end{question}
\begin{solution}
图\ref{fig:at4:example:5:s} 取碰撞点向外为径向正方向，圆速 $v_c=\sqrt{\mu/R}$。碰撞短时忽略引力冲量，分别守恒切向和径向动量，得 $v_t=10v_c/11$、$v_r=-u/11$。
\typeout{ASTROFIG:fig:at4:example:5:p}\fig{fig:at4:example:5:p}{解答过程：陨石粘合后卫星是否撞地}{\begin{tikzpicture}[scale=.85]
\fill(0,0)circle(2pt)node[below]{$F$};\fill(3,0)circle(2pt)node[below]{粘合点};\draw[flow](3,0)--(3,1.5)node[right]{$10v_c/11$};\draw[flow](3,1.5)--(2,1.5)node[above]{$-u/11$};\draw[force](3,0)--(2,1.5);
\end{tikzpicture}}

图\ref{fig:at4:example:5:p} 在碰后飞行段，$h=10Rv_c/11$，$2\varepsilon=(100v_c^2+u^2)/121-2\mu/R$。在表面可达的条件为
\[2\varepsilon+\frac{2\mu}{R_0}-\frac{h^2}{R_0^2}\ge0.
\]
代入并乘121，再令 $x=R/R_0>1$，整理为
\[u^2\ge\frac\mu R(100x^2-242x+142)
=\frac{2\mu}{R}(x-1)(50x-71).
\]
因此 $u_{\min}=\sqrt{\max\{0,\,2\mu(x-1)(50x-71)/R\}}$。当 $1<x\le71/50$，仅切向减速就已使近点降至表面以下，最低径向速率为零；不是所有参数下都存在正阈值。碰撞前后机械能不守恒，守恒的是短时动量与碰后飞行能量两段各自的关系。\end{solution}
\begin{question}\begin{example}[非零入射速度的两体相遇时间]\label{at4:example:6}
两质点相距 $l$，$m_1$ 以速度 $u_0\ge0$ 沿连线朝初始静止的 $m_2$ 运动，只受引力。求相遇时间，给出三种相对能量下的表达式。
\end{example}
\typeout{ASTROFIG:fig:at4:example:6:s}\fig{fig:at4:example:6:s}{题设：非零入射速度的两体相遇时间}{\begin{tikzpicture}[scale=.85]
\fill(0,0)circle(3pt)node[above]{$m_1$};\fill(5,0)circle(3pt)node[above]{$m_2$};\draw[dim](0,-.6)--(5,-.6)node[midway,below]{$l$};\draw[flow](0,.5)--(1.3,.5)node[above]{$u_0$};
\end{tikzpicture}}
\end{question}
\begin{solution}
图\ref{fig:at4:example:6:s} 的初始相对速度为 $-u_0$，相对引力参数 $\mu=G(m_1+m_2)$。能量给 $\dot r=-\sqrt{2\mu/r-k}$，$k=2\mu/l-u_0^2$。
\typeout{ASTROFIG:fig:at4:example:6:p}\fig{fig:at4:example:6:p}{解答过程：非零入射速度的两体相遇时间}{\begin{tikzpicture}[scale=.85]
\fill(0,0)circle(3pt)node[above]{$m_1$};\fill(5,0)circle(3pt)node[above]{$m_2$};\draw[dim](0,-.6)--(5,-.6)node[midway,below]{$l$};\node at(2.5,1){$\dot r<0$};\node at(2.5,-1.7){$k>0,=0,<0$ 分别换元};
\end{tikzpicture}}

图\ref{fig:at4:example:6:p} 对全过程取时间积分 $t=\int_0^l\dd r/\sqrt{2\mu/r-k}$。若 $k>0$，令 $A=2\mu/k$、$r=A\sin^2\theta$，$\theta_0=\arcsin\sqrt{l/A}$，则
\[t=\frac{A}{\sqrt k}(\theta_0-\sin\theta_0\cos\theta_0).
\]
若 $k=0$，直接积 $\sqrt{r/(2\mu)}$，得 $t=2l^{3/2}/(3\sqrt{2\mu})$。若 $k<0$，令 $s=-k>0$、$r=(2\mu/s)\sinh^2\eta$，$\eta_0=\operatorname{arsinh}\sqrt{sl/(2\mu)}$，得
\[t=\frac{2\mu}{s^{3/2}}(\sinh\eta_0\cosh\eta_0-\eta_0).
\]
换元时 $\dd r/\sqrt{2\mu/r-k}$ 分别化为 $2A\sin^2\theta\dd\theta/\sqrt k$ 和 $4\mu\sinh^2\eta\dd\eta/s^{3/2}$，所以两个原函数前的系数不同。$u_0=0$ 时第一式恢复静止相落；大 $u_0$ 时结果趋于 $l/u_0$。尽管高速时能量为正，当前仍是向内的径向碰撞分支。\end{solution}
\begin{question}\begin{example}[垂直初速度下的相互逃逸]\label{at4:example:7}
两质点相距 $l$，$m_1$ 初速度 $v_1$ 沿连线指向 $m_2$，$m_2$ 初速度 $v_2$ 与该连线垂直。只受相互引力，求可以相互分离至无穷远的条件，并检查 $v_2=0$。
\end{example}
\typeout{ASTROFIG:fig:at4:example:7:s}\fig{fig:at4:example:7:s}{题设：垂直初速度下的相互逃逸}{\begin{tikzpicture}[scale=.85]
\fill(0,0)circle(3pt)node[above]{$m_1$};\fill(5,0)circle(3pt)node[above]{$m_2$};\draw[dim](0,-.6)--(5,-.6)node[midway,below]{$l$};\draw[flow](0,.5)--(1.3,.5)node[above]{$v_1$};\draw[flow](5,0)--(5,1.4)node[right]{$v_2$};
\end{tikzpicture}}
\end{question}
\begin{solution}
图\ref{fig:at4:example:7:s} 中相对速度平方为 $v_1^2+v_2^2$，相对角动量为 $h=lv_2$。相对能量非负给
\[v_1^2+v_2^2\ge\frac{2G(m_1+m_2)}l.
\]
\typeout{ASTROFIG:fig:at4:example:7:p}\fig{fig:at4:example:7:p}{解答过程：垂直初速度下的相互逃逸}{\begin{tikzpicture}[scale=.85]
\draw[axis](0,0)--(4,0)node[right]{$v_1$};\draw[axis](0,0)--(0,3)node[above]{$v_2$};\draw[world](0,2)arc(90:0:2);\node at(2.9,2){能量允许区：圆弧外};\node at(2.7,.5){另检验 $v_2>0$};
\end{tikzpicture}}

当 $v_2>0$，图\ref{fig:at4:example:7:p} 的轨道有非零最近距离，点质点模型下绕过后可沿抛物线或双曲线分离，因此上述不等式充分。可写阈值 $v_2\ge\sqrt{2\mu/l-v_1^2}$，但仅在根内非负且保持 $v_2>0$ 时这样表达。

若 $v_2=0,v_1>0$，相对运动径向向内，必先相遇；能量再高也不能直接称作相互逃逸。若 $v_1^2\ge2\mu/l$，任意非零且足够小的 $v_2$ 都给逃逸，所需 $v_2$ 的下确界为零，却不由径向入射的零值实际取得。真实有半径天体还需增加最近距离大于接触距离的条件。\end{solution}
\begin{question}\begin{example}[从地球表面直接逃离太阳系]\label{at4:example:8}
地球在半径 $r_E$ 的圆轨道上绕太阳运动，日心圆速 $V_E$。地球半径 $R_E$、引力参数 $\mu_E$。忽略自转、大气，采用分段二体近似，求直接逃离太阳系的最小地表速率。另说明从地表静止瞬时加速到该速率时，对固定质量试验飞船做的功。
\end{example}
\typeout{ASTROFIG:fig:at4:example:8:s}\fig{fig:at4:example:8:s}{题设：从地球表面直接逃离太阳系}{\begin{tikzpicture}[scale=.85]
\draw[world](0,0)circle(2.3);\fill(0,0)circle(4pt)node[below]{太阳};\fill(2.3,0)circle(3pt)node[below]{地球};\draw[flow](2.3,0)--(2.3,1.4)node[right]{$V_E$};
\end{tikzpicture}}
\end{question}
\begin{solution}
图\ref{fig:at4:example:8:s} 在地球影响区外，日心临界逃逸速率为 $\sqrt2V_E$。为使所需地球系余速最小，应让余速沿地球公转方向，故 $u_E=(\sqrt2-1)V_E$。
\typeout{ASTROFIG:fig:at4:example:8:p}\fig{fig:at4:example:8:p}{解答过程：从地球表面直接逃离太阳系}{\begin{tikzpicture}[scale=.85]
\draw[flow](0,0)--(2,0)node[midway,below]{$V_E$};\draw[flow](2,0)--(3,0)node[midway,above]{$u_E$};\node at(1.5,1){$V_E+u_E=\sqrt2V_E$};\node at(1.5,-1){$v_0^2=u_E^2+2\mu_E/R_E$};
\end{tikzpicture}}

图\ref{fig:at4:example:8:p} 再回到地球系，由地表到影响区外的能量式 $v_0^2-2\mu_E/R_E=u_E^2$，得到
\[v_{0,\min}=\sqrt{\frac{2\mu_E}{R_E}+(\sqrt2-1)^2V_E^2}.
\]
若把地球系中地表初速视为零，固定质量 $m$ 的试验飞船经一个理想瞬时外力加速，做功是 $mv_0^2/2$，而不是两个逃逸能的简单相加。这不等于真实火箭燃料的全部化学能：喷出气体也带走能量，质量亦改变。这里忽略地球自转，若计入则初始速度矢量与发射纬度也会影响结果。\end{solution}
\begin{question}\begin{example}[移动行星旁的末速为何不唯一]\label{at4:example:9}
探测器在距行星中心 $r$ 处的恒星系速率为 $v_0$，与行星恒星系速度 $\bm V$ 垂直。行星引力参数 $\mu_p$，探测器以后能够离开。仅凭这些条件能否唯一确定最终恒星系速度？给出能确定的量及需要补充的条件。
\end{example}
\typeout{ASTROFIG:fig:at4:example:9:s}\fig{fig:at4:example:9:s}{题设：移动行星旁的末速为何不唯一}{\begin{tikzpicture}[scale=.85]
\fill(0,0)circle(4pt)node[below]{行星};\draw[flow](0,0)--(1.5,0)node[below]{$\bm V$};\fill(3,1)circle(2pt)node[right]{探测器};\draw[flow](3,1)--(3,2.2)node[above]{$\bm v_0$};\draw[guide](0,0)--(3,1)node[midway,above]{$r$};
\end{tikzpicture}}
\end{question}
\begin{solution}
图\ref{fig:at4:example:9:s} 转到行星系，当前相对速率平方 $u_0^2=v_0^2+V^2$。行星系能量守恒可确定无穷远余速
\[u_\infty^2=v_0^2+V^2-\frac{2\mu_p}{r}\ge0.
\]
\typeout{ASTROFIG:fig:at4:example:9:p}\fig{fig:at4:example:9:p}{解答过程：移动行星旁的末速为何不唯一}{\begin{tikzpicture}[scale=.85]
\draw[flow](0,0)--(2.3,0)node[midway,below]{$\bm V$};\draw[guide](2.3,0)circle(1.4);\draw[flow](2.3,0)--(3.3,1)node[right]{$\bm u_{\rm out}$};\draw[force](0,0)--(3.3,1)node[above]{$\bm v_{\rm out}$};
\end{tikzpicture}}

图\ref{fig:at4:example:9:p} 回到恒星系，$\bm v_{\rm out}=\bm V+\bm u_{\rm out}$，其大小还含出射方向：
\[v_{\rm out}^2=V^2+u_\infty^2+2Vu_\infty\cos\gamma.
\]
因此仅可给必要的三角不等式范围 $|V-u_\infty|\le v_{\rm out}\le V+u_\infty$，并不能断言整个范围在指定掠近约束下都可实现。必须补充当前相对位置矢量、入射或出射分支，或者等效地给出入射渐近方向、冲量参数和绕行一侧；还要检验不穿入行星。把相对速度直接反转180度只是另一个额外假设，不能从题中两速度垂直推出。\end{solution}
\begin{question}\begin{example}[经过火星一次助推的发射速度条件]\label{at4:example:10}
地球、火星按半径 $r_E<r_M$ 的同向共面圆轨道绕太阳运动，日心引力参数 $\mu_s$；地球半径 $R_E$、参数 $\mu_E$，火星参数 $\mu_M$、最小允许掠近半径 $r_{\min}$。允许等待相位。离地后探测器以日心切向速度 $x$ 顺行，首次向外到达火星轨道，随后无动力掠近火星一次。给出求最小地表发射速率的完整一维条件。
\end{example}
\typeout{ASTROFIG:fig:at4:example:10:s}\fig{fig:at4:example:10:s}{题设：经过火星一次助推的发射速度条件}{\begin{tikzpicture}[scale=.85]
\draw[guide](0,0)circle(1.5);\draw[guide](0,0)circle(3);\fill(0,0)circle(3pt)node[below]{太阳};\fill(1.5,0)circle(2pt)node[right]{地球};\fill(-2.142857,2.099563)circle(2pt)node[above]{火星};\draw[world,domain=0:135.58469,samples=90]plot({2.11764706*cos(\x)/(1+.411764706*cos(\x))},{2.11764706*sin(\x)/(1+.411764706*cos(\x))});
\end{tikzpicture}}
\end{question}
\begin{solution}
图\ref{fig:at4:example:10:s} 的约束明确了出发方向、第一次交会和安全掠近半径。记 $V_E=\sqrt{\mu_s/r_E}$、$V_M=\sqrt{\mu_s/r_M}$。到达火星半径至少需要 $x\ge x_H=V_E\sqrt{2r_M/(r_E+r_M)}$，等号为双切转移。

出发日心能量 $\varepsilon=x^2/2-\mu_s/r_E$，角动量 $h=r_Ex$。到火星半径时
\[v_t=\frac{r_Ex}{r_M},\quad
v_r=\sqrt{x^2-\frac{2\mu_s}{r_E}+\frac{2\mu_s}{r_M}-v_t^2},\quad
U=\sqrt{v_r^2+(v_t-V_M)^2}.
\]
取 $v_r\ge0$ 对应第一次向外到达。相对入射与火星顺行方向的夹角 $\alpha=\operatorname{atan2}(v_r,v_t-V_M)$ 在 $[0,\pi]$。
\typeout{ASTROFIG:fig:at4:example:10:p}\fig{fig:at4:example:10:p}{解答过程：经过火星一次助推的发射速度条件}{\begin{tikzpicture}[scale=.85]
\draw[flow](0,0)--(2.5,0)node[midway,below]{$\bm V_M$};\draw[flow](2.5,0)--(1.5,1.7)node[midway,right]{$\bm U_{\rm in}$};\draw[flow,blue](2.5,0)--(4.4,.6)node[right]{$\bm U_{\rm out}$};\draw[guide](2.5,0)circle(1.97);\node at(2.5,-2.6){允许转角：$0\le\delta\le\delta_{\max}$};
\end{tikzpicture}}

图\ref{fig:at4:example:10:p} 掠近最多可转 $\delta_{\max}=2\arcsin[1/(1+r_{\min}U^2/\mu_M)]$。允许选择绕行一侧时，应尽量把余速朝向火星公转速度；最小剩余夹角为 $\gamma_{\min}=\max(0,\alpha-\delta_{\max})$。于是可实现的最大日心末速平方为
\[W^2(x)=V_M^2+U^2+2V_MU\cos\gamma_{\min}.
\]
太阳逃逸要求 $W^2(x)\ge2\mu_s/r_M$。因此在 $x\ge x_H$ 的可行集合内，取满足该条件的最小 $x_*$，地表所需速率为
\[v_{0,\min}=\sqrt{2\mu_E/R_E+(x_*-V_E)^2}.\]
给出常量后可分段扫描、夹住边界并求根，再回代检查轨道交会与掠近条件；不应未经证明假设整个可行函数单调。若不给最小掠近半径、方向和交会限制，所谓“助推后的唯一最小发射速度”没有唯一的数值定义。本题所给是一套明确近似下可计算的完整条件，未把不受限制的180度反转当成自动可行。\end{solution}
\begin{question}\begin{example}[双曲轨道到达近点的时间]\label{at4:example:11}
引力双曲线正几何半长轴 $a_h$、离心率 $e>1$。探测器在入射支，初半径 $r_0>a_h(e-1)$。求到近点的剩余时间。
\end{example}
\typeout{ASTROFIG:fig:at4:example:11:s}\fig{fig:at4:example:11:s}{题设：双曲轨道到达近点的时间}{\begin{tikzpicture}[scale=.85]
\draw[guide](-4.6,-1.118)--(1,-1.118);\draw[world,domain=-2:2,samples=110]plot({(1.5-cosh(\x))/1.5+1.11803399*1.11803399*sinh(\x)/1.5},{-1.11803399*(1.5-cosh(\x))/1.5+1.11803399*sinh(\x)/1.5});\fill(0,0)circle(2pt)node[above left]{$F$};\draw[flow](-4.2,-1.118)--(-3,-1.118)node[midway,below]{$v_\infty$};
\end{tikzpicture}}
\end{question}
\begin{solution}
  图\ref{fig:at4:example:11:s} 用双曲参数 $r=a_h(e\cosh H-1)$，初态入射选 $H_0<0$。令 $H_* =\operatorname{arcosh}[(r_0/a_h+1)/e]>0$，则 $H_0=-H_*$。
  \typeout{ASTROFIG:fig:at4:example:11:p}\fig{fig:at4:example:11:p}{解答过程：双曲轨道到达近点的时间}{\begin{tikzpicture}[scale=.85]
\draw[guide](-4.6,-1.118)--(1,-1.118);\draw[world,domain=-2:2,samples=110]plot({(1.5-cosh(\x))/1.5+1.11803399*1.11803399*sinh(\x)/1.5},{-1.11803399*(1.5-cosh(\x))/1.5+1.11803399*sinh(\x)/1.5});\fill(0,0)circle(2pt)node[above left]{$F$};\draw[flow](-4.2,-1.118)--(-3,-1.118)node[midway,below]{$v_\infty$};\node at(-2,2){入射 $H_0=-H_*<0$};\node at(-2,1.3){近点 $H=0$};
\end{tikzpicture}}

  图\ref{fig:at4:example:11:p} 在双曲线中心坐标取 $X=a_h\cosh H,Y=b_h\sinh H$，焦点位置相差 $a_he$。扫面积微分大小为 $a_hb_h(e\cosh H-1)\dd H/2$。除以 $h/2$，并用 $h=b_h\sqrt{\mu/a_h}$，得 $\dd t=\sqrt{a_h^3/\mu}(e\cosh H-1)\dd H$。
  从 $-H_*$ 积到零，剩余时间 $t=\sqrt{a_h^3/\mu}(e\sinh H_*-H_*)>0$。初点趋于近点时 $H_*\to0$、时间趋零。这个结果使用实际入射分支，没有把正参数的出射时刻误当成未来近点时间。\end{solution}
\begin{question}\begin{example}[逆行点火使入射天体被捕获]\label{at4:example:12}
粒子余速 $v_\infty>0$，近点 $r_p$，中心参数 $\mu$。在近点作逆行脉冲，求从双曲线变成有界椭圆所需速度减量的下确界，并讨论近点越低是否越有利。
\end{example}
\typeout{ASTROFIG:fig:at4:example:12:s}\fig{fig:at4:example:12:s}{题设：逆行点火使入射天体被捕获}{\begin{tikzpicture}[scale=.85]
\draw[guide](-4.6,-1.118)--(1,-1.118);\draw[world,domain=-2:2,samples=110]plot({(1.5-cosh(\x))/1.5+1.11803399*1.11803399*sinh(\x)/1.5},{-1.11803399*(1.5-cosh(\x))/1.5+1.11803399*sinh(\x)/1.5});\fill(0,0)circle(2pt)node[above left]{$F$};\draw[flow](-4.2,-1.118)--(-3,-1.118)node[midway,below]{$v_\infty$};
\end{tikzpicture}}
\end{question}
\begin{solution}
  图\ref{fig:at4:example:12:s} 点火前 $v_p=\sqrt{v_\infty^2+2\mu/r_p}$。若点火后能量为负，需 $v_+<\sqrt{2\mu/r_p}$，并且保留非零切向速度以免退化为直接碰撞。
  \typeout{ASTROFIG:fig:at4:example:12:p}\fig{fig:at4:example:12:p}{解答过程：逆行点火使入射天体被捕获}{\begin{tikzpicture}[scale=.85]
\draw[guide](-4.6,-1.118)--(1,-1.118);\draw[world,domain=-2:2,samples=110]plot({(1.5-cosh(\x))/1.5+1.11803399*1.11803399*sinh(\x)/1.5},{-1.11803399*(1.5-cosh(\x))/1.5+1.11803399*sinh(\x)/1.5});\fill(0,0)circle(2pt)node[above left]{$F$};\draw[flow](-4.2,-1.118)--(-3,-1.118)node[midway,below]{$v_\infty$};\draw[force](.333,-.373)--(.7,-1.1)node[right]{逆行点火};\node at(-2,2){$\varepsilon_+<0$ 才束缚};
\end{tikzpicture}}

  图\ref{fig:at4:example:12:p} 速度减量的下确界为 $d_* =\sqrt{v_\infty^2+2\mu/r_p}-\sqrt{2\mu/r_p}$，等号只得到抛物线，不是束缚椭圆。真正捕获应取 $d>d_*$，并检验新近点或当前端点是否仍安全。把两根之差有理化，$d_*=v_\infty^2/[\sqrt{v_\infty^2+2\mu/r_p}+\sqrt{2\mu/r_p}]$，可见减小 $r_p$ 会减小这一下界；实际不能越过最小安全半径。下确界与能够实际取得的最小值在这里需要区分。\end{solution}
\begin{question}\begin{example}[给定转向的完整引力助推]\label{at4:example:13}
在行星影响区外，行星日心速率 $V$，探测器行星系余速 $u=V$，入射方向与行星速度垂直。选择绕行一侧，使余速向行星公转方向偏转 $45^\circ$。求日心末速和所需冲量参数，并给出安全近点的检查式。
\end{example}
\typeout{ASTROFIG:fig:at4:example:13:s}\fig{fig:at4:example:13:s}{题设：给定转向的完整引力助推}{\begin{tikzpicture}[scale=.85]
\draw[flow](0,0)--(2,0)node[below]{$\bm V$};\draw[flow](0,0)--(0,2)node[left]{$\bm u_{\rm in}$};\draw[flow,blue](0,0)--(1.414,1.414)node[above]{$\bm u_{\rm out}$};
\end{tikzpicture}}
\end{question}
\begin{solution}
  图\ref{fig:at4:example:13:s} 已给定入射、出射一侧与转角，末余速分量为 $(u/\sqrt2,u/\sqrt2)$。因此日心末速度为 $(V+u/\sqrt2,u/\sqrt2)$。
  \typeout{ASTROFIG:fig:at4:example:13:p}\fig{fig:at4:example:13:p}{解答过程：给定转向的完整引力助推}{\begin{tikzpicture}[scale=.85]
\draw[flow](0,0)--(2,0)node[below]{$V$};\draw[flow,blue](2,0)--(3.414,1.414)node[right]{$u$};\draw[force](0,0)--(3.414,1.414)node[above]{$v_{\rm out}$};
\end{tikzpicture}}

  图\ref{fig:at4:example:13:p} 取 $u=V$ 得 $v_{\rm out}=V\sqrt{2+\sqrt2}$，初日心速率为 $V\sqrt2$，确有增加。散射式 $\tan(\chi/2)=\mu_p/(bu^2)$ 给 $b=\mu_p/[u^2\tan(\pi/8)]$。离心率 $e=1/\sin(\pi/8)$，最近距离 $r_p=\mu_p(e-1)/u^2$，还须 $r_p\ge r_{\min}$。若不满足，所要求45度助推在给定安全约束下不可行，不能仅凭速度三角形断定能实现。\end{solution}
\begin{question}\begin{example}[有限半径双天体的相撞截面]\label{at4:example:14}
两个球形天体质量为 $M_1,M_2$、半径为 $R_1,R_2$，在无穷远相对速度为 $u$，相对冲量参数为 $b$。忽略潮汐形变，求两表面相遇的临界截面。
\end{example}
\typeout{ASTROFIG:fig:at4:example:14:s}\fig{fig:at4:example:14:s}{题设：有限半径双天体的相撞截面}{\begin{tikzpicture}[scale=.85]
\draw[world](0,0)circle(1);\draw[world](1.6,0)circle(.6);\fill(0,0)circle(2pt)node[above]{$M_1$};\fill(1.6,0)circle(2pt)node[above]{$M_2$};\draw[dim](0,-1.4)--(1.6,-1.4)node[midway,below]{$R_1+R_2$};\node at(.8,1.7){接触状态；相对初始参数 $b,u$};
\end{tikzpicture}}
\end{question}
\begin{solution}
  图\ref{fig:at4:example:14:s} 两体都运动时使用相对参数 $\mu=G(M_1+M_2)$，接触距离为 $R_c=R_1+R_2$，不是任一单独半径。相对角动量 $h=bu$，相对能量 $u^2/2$。
  \typeout{ASTROFIG:fig:at4:example:14:p}\fig{fig:at4:example:14:p}{解答过程：有限半径双天体的相撞截面}{\begin{tikzpicture}[scale=.85]
\draw[world](0,0)circle(1);\draw[world](1.6,0)circle(.6);\fill(0,0)circle(2pt)node[above]{$M_1$};\fill(1.6,0)circle(2pt)node[above]{$M_2$};\draw[dim](0,-1.4)--(1.6,-1.4)node[midway,below]{$R_1+R_2$};\node at(.8,1.7){接触状态；相对初始参数 $b,u$};\node at(.8,-2.4){$h=bu,\quad\mu=G(M_1+M_2)$};
\end{tikzpicture}}

  图\ref{fig:at4:example:14:p} 表面接触时径向速度平方 $u^2+2\mu/R_c-b^2u^2/R_c^2$，临界为零，故
  \[\sigma=\pi b_{\max}^2=\pi(R_1+R_2)^2\left[1+\frac{2G(M_1+M_2)}{(R_1+R_2)u^2}\right].\]
  换到质心系后没有必要把两颗球各自的实验室动能再重复加进相对式。若一个质量和尺寸都可忽略，结果恢复粒子撞固定球形天体的聚焦截面。接触后的粘合、反弹或碎裂不由这条截面公式决定。\end{solution}

\section{本讲小结}
逃逸先在正确的相对参考系检验能量，再检验角动量、表面和时间方向。双曲线区分正几何半长轴、冲量参数与最近距离。碰撞短时守恒动量，碰后飞行另用机械能守恒。助推先在行星系转向，再变回恒星系；安全掠近半径限制转角，条件不足的题目应保留参数而不虚构唯一数值。

\section{独立练习}

\begin{question}\begin{drill}[弱偏转的最近距离修正]\label{at4:drill:1}
入射参数 $b,v_\infty$ 满足 $\eta=\mu/(bv_\infty^2)\ll1$，求偏转角和最近距离的一阶近似。
\end{drill}
\typeout{ASTROFIG:fig:at4:drill:1:s}\fig{fig:at4:drill:1:s}{题设：弱偏转的最近距离修正}{\begin{tikzpicture}[scale=.85]
\draw[guide](-4.5,-1.1055)--(4.5,-1.1055);\draw[world,domain=-2:2,samples=100]plot({(.11111111*(10-cosh(\x))+9.94987437*1.1055416*sinh(\x))/10},{(-9.94987437*.11111111*(10-cosh(\x))+1.1055416*sinh(\x))/10});\fill(0,0)circle(2pt)node[above]{$F$};\draw[flow](-4.1,-1.1055)--(-3,-1.1055)node[below]{$v_\infty$};
\end{tikzpicture}}
\end{question}
\hintlink{at4:drill:1}
\begin{question}\begin{drill}[给定余速与安全半径的最大转角]\label{at4:drill:2}
行星参数 $\mu_p$，相对余速 $u$，掠近点不得小于 $r_{\min}$。求最大偏转角和对应冲量参数，并说明能否精确反转相对速度。
\end{drill}
\typeout{ASTROFIG:fig:at4:drill:2:s}\fig{fig:at4:drill:2:s}{题设：给定余速与安全半径的最大转角}{\begin{tikzpicture}[scale=.85]
\draw[guide](-4.6,-1.118)--(1,-1.118);\draw[world,domain=-2:2,samples=110]plot({(1.5-cosh(\x))/1.5+1.11803399*1.11803399*sinh(\x)/1.5},{-1.11803399*(1.5-cosh(\x))/1.5+1.11803399*sinh(\x)/1.5});\fill(0,0)circle(2pt)node[above left]{$F$};\draw[flow](-4.2,-1.118)--(-3,-1.118)node[midway,below]{$u$};\draw[world](0,0)circle(.5);
\end{tikzpicture}}
\end{question}
\hintlink{at4:drill:2}
\begin{question}\begin{drill}[引力助推的能量正负]\label{at4:drill:3}
行星在恒星系速度为 $\bm V$。探测器入、出行星系余速大小均为 $u$，分别与 $\bm V$ 夹角 $\alpha,\beta$。求恒星系单位质量能量变化，并给出加速条件。
\end{drill}
\typeout{ASTROFIG:fig:at4:drill:3:s}\fig{fig:at4:drill:3:s}{题设：引力助推的能量正负}{\begin{tikzpicture}[scale=.85]
\draw[flow](0,0)--(2,0)node[below]{$\bm V$};\draw[flow](0,0)--(1,1.5)node[above]{$\bm u_{\rm in}$};\draw[flow,blue](0,0)--(1.7,.5)node[right]{$\bm u_{\rm out}$};
\end{tikzpicture}}
\end{question}
\hintlink{at4:drill:3}
\begin{question}\begin{drill}[局部速度相同却有不同最近距离]\label{at4:drill:4}
两粒子在相同半径 $r_0$，速率均为 $\sqrt{\mu/r_0}$。一个纯切向，另一个径向向外。比较两条轨道的能量、角动量和最大最小距离。
\end{drill}
\typeout{ASTROFIG:fig:at4:drill:4:s}\fig{fig:at4:drill:4:s}{题设：局部速度相同却有不同最近距离}{\begin{tikzpicture}[scale=.85]
\fill(0,0)circle(2pt)node[below]{$F$};\draw[guide](0,0)circle(2);\fill(2,0)circle(2pt);\draw[flow](2,0)--(2,1.3)node[left]{切向};\draw[flow,blue](2,0)--(3.3,0)node[above]{径向};\draw[dim](0,-.6)--(2,-.6)node[midway,below]{$r_0$};
\end{tikzpicture}}
\end{question}
\hintlink{at4:drill:4}
\begin{question}\begin{drill}[无耗散引力能否自动捕获]\label{at4:drill:5}
粒子从无穷远以非零余速掠过一个静止球形天体，轨道不碰表面。仅靠二体保守引力，能否最终留在椭圆轨道？若能量从其他系统交换，应如何修改判断？
\end{drill}
\typeout{ASTROFIG:fig:at4:drill:5:s}\fig{fig:at4:drill:5:s}{题设：无耗散引力能否自动捕获}{\begin{tikzpicture}[scale=.85]
\draw[guide](-4.6,-1.118)--(1,-1.118);\draw[world,domain=-2:2,samples=110]plot({(1.5-cosh(\x))/1.5+1.11803399*1.11803399*sinh(\x)/1.5},{-1.11803399*(1.5-cosh(\x))/1.5+1.11803399*sinh(\x)/1.5});\fill(0,0)circle(2pt)node[above left]{$F$};\draw[flow](-4.2,-1.118)--(-3,-1.118)node[midway,below]{$v_\infty$};
\end{tikzpicture}}
\end{question}
\hintlink{at4:drill:5}
\begin{question}\begin{drill}[超过给定偏转角的散射截面]\label{at4:drill:6}
固定中心参数 $\mu$，入射速度 $u$，统计偏转角至少为 $\chi_0$ 的轨道，$0<\chi_0<\pi$。求点质量模型下的截面；若中心有吸收半径，再说明怎样排除碰撞。
\end{drill}
\typeout{ASTROFIG:fig:at4:drill:6:s}\fig{fig:at4:drill:6:s}{题设：超过给定偏转角的散射截面}{\begin{tikzpicture}[scale=.85]
\draw[guide](-4.6,-1.118)--(1,-1.118);\draw[world,domain=-2:2,samples=110]plot({(1.5-cosh(\x))/1.5+1.11803399*1.11803399*sinh(\x)/1.5},{-1.11803399*(1.5-cosh(\x))/1.5+1.11803399*sinh(\x)/1.5});\fill(0,0)circle(2pt)node[above left]{$F$};\draw[flow](-4.2,-1.118)--(-3,-1.118)node[midway,below]{$u$};
\end{tikzpicture}}
\end{question}
\hintlink{at4:drill:6}
\begin{question}\begin{drill}[吸收静止尘埃后的减速与耗散]\label{at4:drill:7}
质量 $M$ 的天体以速度 $V$ 运动，短时间完全吸收原先静止的质量 $\dd M$。忽略这段时间外力冲量，求一阶速度变化和机械能损失。
\end{drill}
\typeout{ASTROFIG:fig:at4:drill:7:s}\fig{fig:at4:drill:7:s}{题设：吸收静止尘埃后的减速与耗散}{\begin{tikzpicture}[scale=.85]
\draw[world](0,-.4)rectangle(1.5,.4);\node at(.75,0){$M$};\draw[flow](0,1)--(1.5,1)node[above]{$V$};\fill(3,0)circle(2pt)node[above]{$\dd M$ 静止};
\end{tikzpicture}}
\end{question}
\hintlink{at4:drill:7}
\begin{question}\begin{drill}[恒星系助推能量的上界为何不能随意取得]\label{at4:drill:8}
行星日心速度大小 $V$，探测器相对余速 $u$。证明任一次无动力掠近的单位质量日心能量变化满足 $|\Delta\varepsilon|\le2Vu$；若转角最多 $\delta_{\max}$，给出更紧的方向无关上界。
\end{drill}
\typeout{ASTROFIG:fig:at4:drill:8:s}\fig{fig:at4:drill:8:s}{题设：恒星系助推能量的上界为何不能随意取得}{\begin{tikzpicture}[scale=.85]
\draw[guide](0,0)circle(1.6);\draw[flow](0,0)--(1.6,0)node[right]{$\bm u_{\rm in}$};\draw[flow,blue](0,0)--(-.8,1.386)node[above]{$\bm u_{\rm out}$};
\end{tikzpicture}}
\end{question}
\hintlink{at4:drill:8}

\LectureAnswerSection\subsection*{两级提示}
\phantomsection\label{hint:at4:drill:1}
\LectureHint{at4:drill:1}{先把全部表达式写成无量纲小参数。}{最近距离的差与相对差分开比较。}
\phantomsection\label{hint:at4:drill:2}
\LectureHint{at4:drill:2}{转角随近点距离增大而减小。}{安全半径为正时检查反正弦的严格范围。}
\phantomsection\label{hint:at4:drill:3}
\LectureHint{at4:drill:3}{将两个恒星系动能展开，相同平方项相消。}{判断沿行星速度的投影是否增加。}
\phantomsection\label{hint:at4:drill:4}
\LectureHint{at4:drill:4}{相同速率和位置只保证相同能量。}{另算角动量，零值必须单独处理。}
\phantomsection\label{hint:at4:drill:5}
\LectureHint{at4:drill:5}{比较入射正能量与椭圆负能量。}{若能量交换存在，明确选取的系统边界。}
\phantomsection\label{hint:at4:drill:6}
\LectureHint{at4:drill:6}{偏转角阈值转成冲量参数上限。}{若只统计掠过，减去会撞入表面的中心截面。}
\phantomsection\label{hint:at4:drill:7}
\LectureHint{at4:drill:7}{先用短时动量关系。}{动能求微分时质量和速度都在变化。}
\phantomsection\label{hint:at4:drill:8}
\LectureHint{at4:drill:8}{能量差是行星速度与余速变化的点积。}{两等长速度矢量之差是速度圆的一条弦。}
\Needspace{8\baselineskip}\subsection*{详解}
\begin{drillsolution}{at4:drill:1}
图\ref{fig:at4:drill:1:s} 的精确关系是 $\chi=2\arctan\eta$，故 $\chi=2\eta+O(\eta^3)$。最近距离由 $r_p=(\mu/v_\infty^2)(\sqrt{1+\eta^{-2}}-1)$ 得
\[r_p=b\big(\sqrt{1+\eta^2}-\eta\big)=b(1-\eta)+O(b\eta^2).
\]
\typeout{ASTROFIG:fig:at4:drill:1:p}\fig{fig:at4:drill:1:p}{解答过程：弱偏转的最近距离修正}{\begin{tikzpicture}[scale=.85]
\draw[guide](-4.5,-1.1055)--(4.5,-1.1055);\draw[world,domain=-2:2,samples=100]plot({(.11111111*(10-cosh(\x))+9.94987437*1.1055416*sinh(\x))/10},{(-9.94987437*.11111111*(10-cosh(\x))+1.1055416*sinh(\x))/10});\fill(0,0)circle(2pt)node[above]{$F$};\draw[flow](-4.1,-1.1055)--(-3,-1.1055)node[below]{$v_\infty$};\node at(0,1){$r_p\simeq b-\mu/v_\infty^2$};
\end{tikzpicture}}

图\ref{fig:at4:drill:1:p} 中最近距离比冲量参数略小，差约 $\mu/v_\infty^2$。虽然绝对修正形式不含 $b$，其相对修正 $\eta$ 会随 $b$ 增大而变小；只有在 $\eta\ll1$ 下这个展开才成立。不能把角度的一阶近似与任意强引力近点同时使用。\end{drillsolution}
\begin{drillsolution}{at4:drill:2}
图\ref{fig:at4:drill:2:s} 有 $e=1+r_pu^2/\mu_p$，$\chi=2\arcsin(1/e)$ 随 $r_p$ 增大而减小。因此 图\ref{fig:at4:drill:2:p} 的最小允许近点给
\[\chi_{\max}=2\arcsin\frac1{1+r_{\min}u^2/\mu_p},\qquad
b_{\min}=r_{\min}\sqrt{1+\frac{2\mu_p}{r_{\min}u^2}}.
\]
\typeout{ASTROFIG:fig:at4:drill:2:p}\fig{fig:at4:drill:2:p}{解答过程：给定余速与安全半径的最大转角}{\begin{tikzpicture}[scale=.85]
\draw[guide](-4.6,-1.118)--(1,-1.118);\draw[world,domain=-2:2,samples=110]plot({(1.5-cosh(\x))/1.5+1.11803399*1.11803399*sinh(\x)/1.5},{-1.11803399*(1.5-cosh(\x))/1.5+1.11803399*sinh(\x)/1.5});\fill(0,0)circle(2pt)node[above left]{$F$};\draw[flow](-4.2,-1.118)--(-3,-1.118)node[midway,below]{$u$};\node at(-2,2){$r_p=r_{\min}$};\node at(-2,1.3){$\chi_{\max}<\pi$};
\end{tikzpicture}}

若 $r_{\min}>0,u>0$，反正弦的自变量严格小于一，转角严格小于 $\pi$，故不能精确180度反转。小余速可以使转角接近180度，却不能忽略对应的长接近时间、行星影响区和其他天体扰动。\end{drillsolution}
\begin{drillsolution}{at4:drill:3}
图\ref{fig:at4:drill:3:s} 的两种恒星系速度为 $\bm V+\bm u_{\rm in,out}$。近似在相同日心位置比较时，恒星势能相同，行星系的 $u^2/2$ 也相同，展开两动能之差得到
\[\Delta\varepsilon=\frac12\big(|\bm V+\bm u_{\rm out}|^2-|\bm V+\bm u_{\rm in}|^2\big)
=Vu(\cos\beta-\cos\alpha).
\]
\typeout{ASTROFIG:fig:at4:drill:3:p}\fig{fig:at4:drill:3:p}{解答过程：引力助推的能量正负}{\begin{tikzpicture}[scale=.85]
\draw[axis](0,0)--(4,0)node[right]{沿 $\bm V$ 的分量};\draw[flow](0,1)--(1.3,1)node[right]{$u\cos\alpha$};\draw[flow,blue](0,2)--(3,2)node[right]{$u\cos\beta$};
\end{tikzpicture}}

图\ref{fig:at4:drill:3:p} 的出射余速更朝向行星速度，即 $\cos\beta>\cos\alpha$，才给日心能量增加；反向绕过可能减速。由余弦范围有绝对上界 $|\Delta\varepsilon|\le2Vu$，但要达到该上界需180度转向，有限安全半径通常不允许。\end{drillsolution}
\begin{drillsolution}{at4:drill:4}
图\ref{fig:at4:drill:4:s} 两初态能量均为 $-\mu/(2r_0)$，所以有相同的有符号半长轴尺度 $r_0$。切向初态 $h=\sqrt{\mu r_0}$，得到 $e=0$，轨道为半径 $r_0$ 的圆。
\typeout{ASTROFIG:fig:at4:drill:4:p}\fig{fig:at4:drill:4:p}{解答过程：局部速度相同却有不同最近距离}{\begin{tikzpicture}[scale=.85]
\draw[world](0,0)circle(1.7);\draw[guide](0,0)--(3.4,0);\fill(3.4,0)circle(2pt)node[right]{$2r_0$};\node at(0,2.3){圆：$r=r_0$};\node at(2.4,-.8){径向：$r_{\max}=2r_0$};
\end{tikzpicture}}

图\ref{fig:at4:drill:4:p} 径向初态 $h=0$，是退化径向轨道，最高距离由 $-\mu/r_{\max}=-\mu/(2r_0)$ 给 $2r_0$，随后掉回中心或实际表面。可把点质点数学极限下的最近距离记零，但遇到实体表面就应停止延伸。相同机械能只固定轨道大小尺度，角动量还决定形状和可达边界。\end{drillsolution}
\begin{drillsolution}{at4:drill:5}
  图\ref{fig:at4:drill:5:s} 初始单位质量相对能量 $\varepsilon=v_\infty^2/2>0$，保守固定中心场中它始终不变。图\ref{fig:at4:drill:5:p} 的轨道虽然改变方向，远处速率仍为 $v_\infty$，故不能自动成为负能量椭圆。
  \typeout{ASTROFIG:fig:at4:drill:5:p}\fig{fig:at4:drill:5:p}{解答过程：无耗散引力能否自动捕获}{\begin{tikzpicture}[scale=.85]
\draw[flow](0,0)--(4,0);\node[above]at(0,0){$\varepsilon>0$ 入射};\node[above]at(4,0){$\varepsilon>0$ 出射};\node at(2,-1){无耗散二体不能自动转为负能量};
\end{tikzpicture}}

  若增加阻力、逆行点火或非弹性碰撞，探测器自身的机械能可以下降；若行星本身移动并与另一天体共同作用，应在完整系统中追踪能量交换。总能量守恒没有失效，但“探测器相对某一固定中心的二体能量”未必仍守恒。不同情形改变的是可使用的守恒量与系统边界。\end{drillsolution}
\begin{drillsolution}{at4:drill:6}
  图\ref{fig:at4:drill:6:s} 由 $\tan(\chi/2)=\mu/(bu^2)$，偏转角随 $b$ 增大而减小。因此 $\chi\ge\chi_0$ 要求 $0\le b\le b_0=\mu/[u^2\tan(\chi_0/2)]$。
  \typeout{ASTROFIG:fig:at4:drill:6:p}\fig{fig:at4:drill:6:p}{解答过程：超过给定偏转角的散射截面}{\begin{tikzpicture}[scale=.85]
\draw[world](0,0)circle(1.8);\draw[fill=gray!20](0,0)circle(.9);\node at(0,0){碰撞区};\node at(1.2,1){散射环};\node at(2.8,-.4){外半径 $b_0$};
\end{tikzpicture}}

  图\ref{fig:at4:drill:6:p} 点质量几何截面为 $\sigma=\pi b_0^2$。若有吸收半径 $R$，非碰撞轨道还需 $b>b_c=R\sqrt{1+2\mu/(Ru^2)}$。同时满足的大角无碰撞截面为 $\pi\max(0,b_0^2-b_c^2)$；边界是否允许相切不影响连续截面的面积值。不能把会撞入中心体的强偏转轨道同时计入“无碰撞散射”。\end{drillsolution}
\begin{drillsolution}{at4:drill:7}
  图\ref{fig:at4:drill:7:s} 初动量 $MV$，末动量 $(M+\dd M)(V+\dd V)$。保留一阶得 $M\dd V+V\dd M=0$，故 $\dd V=-V\dd M/M$。
  \typeout{ASTROFIG:fig:at4:drill:7:p}\fig{fig:at4:drill:7:p}{解答过程：吸收静止尘埃后的减速与耗散}{\begin{tikzpicture}[scale=.85]
\node[draw]at(0,1){$M\dd V=-V\dd M$};\node[draw]at(0,-.5){$\dd K=MV\dd V+V^2\dd M/2$};
\end{tikzpicture}}

  图\ref{fig:at4:drill:7:p} 末动能展开为 $MV^2/2+MV\dd V+V^2\dd M/2$，代入上式后，动能变化为 $-V^2\dd M/2$。因此不可逆机械能损失为 $V^2\dd M/2$，转入内能、形变等。仅计算 $MV\dd V$ 会多算一倍损失，漏掉质量增加本身携带的动能项。若把远处引力聚焦也包含在系统内，还须一致追踪势能与内能。\end{drillsolution}
\begin{drillsolution}{at4:drill:8}
  图\ref{fig:at4:drill:8:s} 两余速大小相同，$\Delta\varepsilon=\bm V\cdot(\bm u_{\rm out}-\bm u_{\rm in})$，因此柯西不等式给 $|\Delta\varepsilon|\le V|\Delta\bm u|$。
  \typeout{ASTROFIG:fig:at4:drill:8:p}\fig{fig:at4:drill:8:p}{解答过程：恒星系助推能量的上界为何不能随意取得}{\begin{tikzpicture}[scale=.85]
\draw[guide](0,0)circle(1.6);\draw[flow](0,0)--(1.6,0);\draw[flow,blue](0,0)--(-.8,1.386);\draw[force](1.6,0)--(-.8,1.386)node[midway,above]{$|\Delta\bm u|$};
\end{tikzpicture}}

  图\ref{fig:at4:drill:8:p} 速度圆上的弦长 $|\Delta\bm u|=2u\sin(\delta/2)$，所以
  \[|\Delta\varepsilon|\le2Vu\sin(\delta_{\max}/2)\le2Vu.
  \]
  前一个界要取得还需速度变化方向与 $\bm V$ 平行；若入射方向已经固定，也未必能达到。安全近点给 $\delta_{\max}<\pi$ 时，后一较宽上界严格不可达。上界用于排除不可能方案，不应在未检查几何条件时直接当作实际能量收益。\end{drillsolution}

\end{document}
